\icode{E-MM-CA-L3-261014-003}
根据量子理论，光子具有动量。光子的动量等于光子的能量除以光速，即\(p = \frac{E}{c}\)。光照射到物体表面并被反射时，会对物体产生压强，这就是“光压”。光压是光的粒子性的典型表现。光压的产生机理如同气体压强：由大量气体分子与器壁的频繁碰撞产生了持续均匀的压力，器壁在单位面积上受到的压力就是气体的压强。
\begin{enumerate}
	\item 激光器发出的一束激光的功率为\(P\)，光束的横截面积为\(S\)。当该激光束垂直照射在物体表面时，试计算单位时间内到达物体表面的光子的总动量。
	\item 若该激光束被物体表面完全反射，试求出其在物体表面引起的光压表达式。
\end{enumerate}
\begin{answer}
    
\end{answer}